%% Wikipedia "Health", Abschnitt "Determinants" (ohne Unterabschnitte),
%% Fassung oldid=1376907986 vom 26.9.2026, abgerufen am 28.10.2026.
%% Lizenz CC BY-SA 4.0; Text unveraendert, nur Bilder und Formatierung angepasst.
%% Erzeugt aus der MediaWiki-API; wird in handouts/LE04.tex und handouts/en/LE04.tex eingebunden.
The context in which an individual lives can impact their health status and quality of life. Many assert health is maintained and improved not only through the advancement and application of health science, but also through the efforts and lifestyle choices of the individual and society. According to the World Health Organization, the main determinants of health include the social and economic environment, the physical environment, and the person's individual characteristics and behaviors.\textsuperscript{[12]}


More specifically, key factors that have been found to influence whether people are healthy or unhealthy include the following:\textsuperscript{[12]}\textsuperscript{[13]}\textsuperscript{[14]}

\begin{itemize}\setlength{\itemsep}{0pt}
\item Education and literacy
\item Employment/working conditions
\item Income and social status
\item Physical environments
\item Social environments
\item Social support networks
\item Biology and genetics
\item Culture
\item Gender
\item Health care services
\item Healthy child development
\item Personal health practices and coping skills
\end{itemize}

Opinion polls shows the public tend to perceive access to health care, health education, and neighborhood environment as most important determinants of good health.\textsuperscript{[15]} An increasing number of studies and reports from different organizations and contexts examine the linkages between health and different factors, including lifestyles, environments, health care organization and health policy, one specific health policy brought into many countries in recent years was the introduction of the sugar tax. Beverage taxes came into light with increasing concerns about obesity, particularly among youth. Sugar-sweetened beverages have become a target of anti-obesity initiatives with increasing evidence of their link to obesity.\textsuperscript{[16]}—such as the 1974 Lalonde report from Canada;\textsuperscript{[14]} the Alameda County Study in California;\textsuperscript{[17]} and the series of World Health Reports of the World Health Organization, which focuses on global health issues including access to health care and improving public health outcomes, especially in developing countries.\textsuperscript{[18]}

The concept of the "health field," as distinct from medical care, emerged from the Lalonde report from Canada. The report identified three interdependent fields as key determinants of an individual's health. These are:\textsuperscript{[14]}

\begin{itemize}\setlength{\itemsep}{0pt}
\item Biomedical: all aspects of health, physical and mental, developed within the human body as influenced by genetic make-up.
\item Environmental: all matters related to health external to the human body and over which the individual has little or no control;
\item Lifestyle: the aggregation of personal decisions (i.e., over which the individual has control) that can be said to contribute to, or cause, illness or death;
\end{itemize}

The maintenance and promotion of health is achieved through different combination of physical, mental, and social well-being—a combination sometimes referred to as the "health triangle."\textsuperscript{[19]} The WHO's 1986 Ottawa Charter for Health Promotion further stated that health is not just a state, but also "a resource for everyday life, not the objective of living. Health is a positive concept emphasizing social and personal resources, as well as physical capacities."\textsuperscript{[20]}

Focusing more on lifestyle issues and their relationships with functional health, data from the Alameda County Study suggested that people can improve their health via exercise, enough sleep, spending time in nature, maintaining a healthy body weight, limiting alcohol use, and avoiding smoking.\textsuperscript{[21]} Health and illness can co-exist, as even people with multiple chronic diseases or terminal illnesses can consider themselves healthy.\textsuperscript{[22]}

\begin{quote}\itshape If you want to learn about the health of a population, look at the air they breathe, the water they drink, and the places where they live.\textsuperscript{[23]}\textsuperscript{[24]}
\\\upshape --- Hippocrates, the Father of Medicine, 5th century BC\end{quote}

The environment is often cited as an important factor influencing the health status of individuals. This includes characteristics of the natural environment, the built environment and the social environment. Factors such as clean water and air, adequate housing, and safe communities and roads all have been found to contribute to good health, especially to the health of infants and children.\textsuperscript{[12]}\textsuperscript{[25]} Some studies have shown that a lack of neighborhood recreational spaces including natural environment leads to lower levels of personal satisfaction and higher levels of obesity, linked to lower overall health and well-being.\textsuperscript{[26]} It has been demonstrated that increased time spent in natural environments is associated with improved self-reported health,\textsuperscript{[27]} suggesting that the positive health benefits of natural space in urban neighborhoods should be taken into account in public policy and land use.

Genetics, or inherited traits from parents, also play a role in determining the health status of individuals and populations. This can encompass both the predisposition to certain diseases and health conditions, as well as the habits and behaviors individuals develop through the lifestyle of their families. For example, genetics may play a role in the manner in which people cope with stress, either mental, emotional or physical. For example, obesity is a significant problem in the United States that contributes to poor mental health and causes stress in the lives of many people.\textsuperscript{[28]} One difficulty is the issue raised by the debate over the relative strengths of genetics and other factors; interactions between genetics and environment may be of particular importance.

\medskip\textbf{References (numbering as in the article)}
{\footnotesize\begin{description}\setlength{\itemsep}{0pt}
\item[{[12]}] World Health Organization. The determinants of health. Geneva. Accessed 12 May 2011.
\item[{[13]}] Public Health Agency of Canada. What Determines Health? Ottawa. Accessed 12 May 2011.
\item[{[14]}] Lalonde, Marc (1974). "A New Perspective on the Health of Canadians." Ottawa: Minister of Supply and Services. Archived 28 October 2014 at the Wayback Machine
\item[{[15]}] Abdalla, Salma M; Hernandez, Mark; Fazaludeen Koya, Shaffi; Rosenberg, Samuel B; Robbins, Grace; Magana, Laura; Nsoesie, Elaine Okanyene; Sabin, Lora; Galea, Sandro (2022). "What matters for health? Public views from eight countries". BMJ Global Health. 7 (6) e008858. doi:10.1136/bmjgh-2022-008858. ISSN 2059-7908. PMC 9204457. PMID 35705225. Retrieved 2026-08-25.
\item[{[16]}] Andreyeva, Tatiana; Chaloupka, Frank J.; Brownell, Kelly D. (2011). "Estimating the potential of taxes on sugar-sweetened beverages to reduce consumption and generate revenue". Preventive Medicine. 52 (6): 413–16. doi:10.1016/j.ypmed.2011.03.013. PMID 21443899.
\item[{[17]}] Housman, Jeff; Dorman, Steve (September–October 2005). "The Alameda County Study: A Systematic, Chronological Review" (PDF). American Journal of Health Education. 36 (5): 302–08. doi:10.1080/19325037.2005.10608200. ISSN 1055-6699. S2CID 39133965. ERIC document number EJ792845. Retrieved 27 December 2011.
\item[{[18]}] World Health Organization. The world health report. Geneva.
\item[{[19]}] Georgia State University. 1998. Health Triangle Slides Archived 11 August 2019 at the Wayback Machine
\item[{[20]}] World Health Organization. The Ottawa Charter for Health Promotion. Adopted at the First International Conference on Health Promotion, Ottawa, 21 November 1986 – WHO/HPR/HEP/95.1.
\item[{[21]}] Housman \& Dorman 2005, pp. 303–04. "The linear model supported previous findings, including regular exercise, limited alcohol consumption, abstinence from smoking, sleeping 7–8 hours a night, and maintenance of a healthy weight play an important role in promoting longevity and delaying illness and death." Citing Wingard DL, Berkman LF, Brand RJ (1982). "A multivariate analysis of health-related practices: a nine-year mortality follow-up of the Alameda County Study". Am J Epidemiol. 116 (5): 765–75. doi:10.1093/oxfordjournals.aje.a113466. PMID 7148802.
\item[{[22]}] Jadad, A.R. (2013). "On Living a Long, Healthy, and Happy Life, Full of Love, and with no Regrets, until Our Last Breath". Verhaltenstherapie. 23 (4): 287–89. doi:10.1159/000357490.
\item[{[23]}] "Environment \& Health – Kids Environment Kids Health NIEHS".
\item[{[24]}] "Environmental Health – Geological Survey Ireland".
\item[{[25]}] UNESCO. The UN World Water Development Report: Facts and Figures – Meeting basic needs. Accessed 12 May 2011.
\item[{[26]}] Björk J, Albin M, Grahn P, Jacobsson H, Ardö J, Wadbro J, Ostergren PO (2008). "Recreational Values of the Natural Environment in Relation to Neighborhood Satisfaction, Physical Activity, Obesity and Well being" (PDF). Journal of Epidemiology \& Community Health. 62 (4): e2. doi:10.1136/jech.2007.062414. PMID 18365329. S2CID 13859588.
\item[{[27]}] White, Mathew P.; Alcock, Ian; Grellier, James; Wheeler, Benedict W.; Hartig, Terry; Warber, Sara L.; Depledge, Michael H.; Fleming, Lora E. (13 June 2019). "Spending at least 120 minutes a week in nature is associated with good health and wellbeing". Scientific Reports. 9 (1): 7730. Bibcode:2019NatSR...9.7730W. doi:10.1038/s41598-019-44097-3. ISSN 2045-2322. PMC 6565732. PMID 31197192.
\item[{[28]}] Kantola, Jussi Ilari; Barath, Tibor; Nazir, Salman; Andre, Terence (2017). Advances in Human Factors, Business Management, Training and Education. Advances in Intelligent Systems and Computing. Vol. 498. doi:10.1007/978-3-319-42070-7. ISBN 978-3-319-42069-1.
\end{description}}
